% !Mode:: "TeX:UTF-8"
%% 请使用 XeLaTeX 编译本文. 适用于 TeX Live 2015 . 本模板的更新下载地址: http://aff.whu.edu.cn/huangzh/
\documentclass[forlib]{WHUPhd}  % 选项 [forprint]: 交付打印时, 加上此选项, 以消除链接文字之彩色, 避免打印偏淡.
                                      % 选项 [forlib]: 提交给图书馆的电子版, 需要加上选项 forlib, 以消除空白页和彩色链接.
% \usepackage{CJKfntef}
%%%%%%%%%%%%%%%%%%%%%=== 参考文献=== %%%%%%%%%%%%%%%%%%%%%
% \bibliographystyle{abbrv}        % 参考文献样式,  plain,unsrt,alpha,abbrv 等等
% \bibliographystyle{bibdata/gbt7714-2005-plain}
\bibliographystyle{bibdata/gbt7714-numerical}
% \bibliographystyle{bibdata/gbt7714-author-year}
% \bibliographystyle{bibdata/ustcthesis-numerical}
\usepackage[numbers]{bibdata/gbt7714}
\usepackage{natbib}

\usepackage{bm}

\usepackage{multirow}
\usepackage{subfigure}
\usepackage{arydshln}
\usepackage{color}

\usepackage{booktabs}


\usepackage{algorithm}
\usepackage{algorithmic}

\renewcommand{\algorithmicrequire}{ \textbf{输入:}} %Use Input in the format of Algorithm
\renewcommand{\algorithmicensure}{ \textbf{输出:}}



% \usepackage[linesnumbered,ruled]{algorithm2e}



\definecolor{mred}{RGB}{255,114,89}
\definecolor{mblue}{RGB}{102, 102, 255}


\definecolor{blue-1}{RGB}{31,139,191}
\definecolor{red-1}{RGB}{219,93,93}



\newcommand{\myparagraph}[1]{{\vspace{6pt}\noindent\heiti #1\vspace{6pt}}}

\newcommand{\note}[1]{\textcolor{red}{【备注:#1】}}


\setlength{\bibsep}{0.5ex}
%%%%%%%%%%%%%%%%%%%%%%%%%%%%%%%%%%%%%%%%%%%%%%%%%%%%

















\begin{document}
%%%%%%%-------------------------------------------------

\fenleihao{TP309}  % 分类号:《中国图书资料分类法》的类号. 必填. 要根据自己的学科方向填写!!
\miji{公开}                % 密级
\UDC{}               %《国际十进制分类法UDC》的类号. 选填.
\bianhao{10486}  % 学校编号, 10486是武汉大学的编号. 不用改动.

\title{博士论文题目-封面}
\titleh{博士论文题目-header}
\Etitle{Phd thesis title} % 英文题目
\author{xxx}
\Eauthor{Name}            %作者英文名
\Csupervisor{xxx\quad 教授}        %指导教师中文名、职称
\Esupervisor{Prof.~xxx}     %指导教师英文名、职称
\Cmajor{xxx}                   %  学科、专业名称
\Emajor{xxx}% 专业英文名
\Cspeciality{xxxx}                     % 研究方向
\Especiality{xxxx}   % 研究方向
\Schoolname{xxx} %学院英文名. 不确定的话, 请看一下自己学院的网页上是怎么写的. 别搞错了!
\date{二〇二一年五月一日}               %  要注意和英文日期一致!!
\Edate{October 1, 2021}                 % 英文封面日期




%-----------------------------------------------------------------------------
\pdfbookmark[0]{封面}{title}         % 封面页加到 pdf 书签

\maketitle

%-----------------------------------------------------------------------------
% !Mode:: "TeX:UTF-8"

%%% 说明: 此部分需要自己填写的内容:  论文创新点.

%%% ``郑重申明'' 和``论文使用授权协议书'' 分别在 statement.tex 和 protocol.tex 中, 无需改动.

%%%%%%%%%%%%%%%%%%%%%%%%%%%%%
%%% -------------  英文封面 (无需改动)-------------   %%%
%%%%%%%%%%%%%%%%%%%%%%%%%%%%%
\thispagestyle{empty}
\renewcommand{\baselinestretch}{1.5}  %下文的行距
% \vspace*{0.5cm}


\begin{center}{\zihao{3} 
A Dissertation Submitted to \\
Academic Degrees Evaluation Committee of Wuhan University for \\
The Degree of Doctor of Philosophy in Computer Science
\par}\end{center}

\vspace*{2cm}

\begin{center}{\zihao{2}\bf \the\Etitle \par}\end{center}

\vfill

\begin{center}
\zihao{4}
\begin{tabular}{ r l }
 Candidate:      &  {\sc \the\Eauthor}      \\
 Supervisor:     &  {\sc \the\Esupervisor}   \\
 Major:          & \the\Emajor  \\
 Speciality:     & \the\Especiality
\end{tabular}

\vspace*{2cm}
\begin{center}
  \iflib % 向图书馆提交电子文档, 使用黑白校徽.
  \includegraphics[height=4cm]{whu.eps}       %%  黑白的. 很小, 只有 10k.
  \else
  \ifprint % 文档打印, 使用黑白校徽.
  \includegraphics[height=4cm]{whu.eps}       %%  黑白的.
  \else
  \includegraphics[height=4cm]{whulogo.eps} %%  彩色的.
  \fi
  \fi
\end{center}


\zihao{-2}
\the\Schoolname\\
{\sc Wuhan University}

\vspace*{1.0cm}

\the\Edate

\end{center}
%%%%%%%--判断是否需要空白页-----------------------------
  \iflib
  \else
  \newpage
  \thispagestyle{empty}
  \cleardoublepage
  \fi
%%%%%%%-------------------------------------------------
{\pagestyle{empty}
%%%--- 加入``郑重声明'' --- %%%%%%%%%%%%%%%%%
% !Mode:: "TeX:UTF-8"
%%-----------------------论文原创性声明-----------------------------------------%
\newpage
\vspace*{20pt}
\begin{center}{\ziju{0.8}\zihao{-2}\heiti 论文原创性声明}\end{center}
\par\vspace*{30pt}
\renewcommand{\baselinestretch}{2}
{\zihao{4} \songti %



本人郑重声明: 所呈交的学位论文, 是本人在导师指导下, 独立进行研究工作所取得的研究成果.
除文中已经标明引用的内容外, 本论文不包含任何其他个人或集体已发表或撰写的研究成果.
对本文的研究做出贡献的个人和集体, 均已在文中以明确方式标明.
本声明的法律结果由本人承担. 

\vskip2cm

\hspace*{4cm}学位论文作者(签名): \hspace{4cm} \hfill \\[1cm]
\hspace*{10cm}年 \hfill  月 \hfill 日\hspace{1cm}\hfill\par}

%%%%%%%--判断是否需要空白页-----------------------------
  \iflib
  \else
  \newpage
  \thispagestyle{empty}
  \cleardoublepage
  \fi
%%%%%%%-------------------------------------------------
\renewcommand{\baselinestretch}{1.6}
\small\normalsize
%%%%%%%%%%%%%%%%%
%%% ---加入``武汉大学学位论文使用授权协议书'' ---  %%%%%%
% !Mode:: "TeX:UTF-8"

%%%%%---武汉大学学位论文使用授权协议书---%%%%%%%%%%%%
%%%%%%%%%%%%%%%%%%%%%%%%%%%%%%%%%%%
%%%%%%%%%%%%%%%%%%%%%%%%%%%%%%%%%%%
\newpage\vspace*{20pt}\thispagestyle{empty}
\begin{center}{\zihao{-2}\heiti 武汉大学学位论文使用授权协议书}\end{center}
\par\vspace*{30pt}

本学位论文作者愿意遵守武汉大学关于保存、使用学位论文的管理办法及规定,
即：学校有权保存学位论文的印刷本和电子版, 并提供文献检索与阅览服务；
学校可以采用影印、缩印、数字化或其它复制手段保存论文；
在以教学与科研服务为目的前提下, 学校可以在校园网内公布部分及全部内容.
\begin{enumerate}[1、]
  \item  在本论文提交当年, 同意在校园网内以及中国高等教育文献保障系统(CALIS)、高校学位论文系统提供查询及前十六页浏览服务.
  \item  在本论文提交~$\Box$~当年/~$\Box$~一年/~$\Box$~两年
            /~$\Box$~三年以后，同意在校园网内允许读者在线浏览并下载全文，学校可以为存在馆际合作关系的兄弟高校用户提供文献传递服务和交换服务.(保密论文解密后遵守此规定)
\end{enumerate}

\vskip 15mm

论文作者(签名)：\raisebox{-1ex}{\underline{\makebox[5cm][c]{}}}
\vskip2em
				          				
学\qquad\qquad\quad 号：\raisebox{-1ex}{\underline{\makebox[5cm][c]{}}}
\vskip2em	
					
学\qquad\qquad\quad 院：\raisebox{-1ex}{\underline{\makebox[5cm][c]{}}}					

\vskip  2cm
\begin{flushright}
 日期：\hskip2cm 年\hskip1.2cm 月\hskip1.2cm 日
\end{flushright}

%%%%%%%%%%%%%%%%%%%%%%%%%%%%%%%%%%%%%%%
%%%%%%%--判断是否需要空白页-----------------------------
  \iflib
  \else
  \newpage
  \thispagestyle{empty}
  \cleardoublepage
  \fi
%%%%%%%-------------------------------------------------
%%%%%%%%%%%%%%%%%%
%%%%%%%%%%%%%%%%%%%%%%%%%%%%%%%
}




%%%%%%%%%%%%%%%%%%%%%%%%%%%%%%%
%%% ------------------- 论文创新点----------------------- %%%
%%%%%%%%%%%%%%%%%%%%%%%%%%%%%%%

\newpage\vspace*{20pt}\thispagestyle{empty}
\begin{center}{\zihao{-2}\heiti 论文创新点}\end{center}
\par\vspace*{30pt}
\baselineskip=20pt



本文的主要创新点包括以下几个部分：
xxxxxxxxxxxxxxxxxxxxx
x

x
x
x
x
x

x
x
x
xxxxxxxxxxxxxx























%%%%%%%%%%%%%%%%%%%%%%%%%%%%%%%
%%%%%%%--请勿删除以下内容 -------------------------------%
%%%%%%%--判断是否需要空白页-----------------------------%
  \iflib
  \let\cleardoublepage\clearpage
  \else
  \newpage
  \thispagestyle{empty}
  \cleardoublepage
  \fi
%%%%%%%---------------------------------------------------%



    % 加入英文封面, 论文创新点, 等等目录之前的部分.
\frontmatter
\pagenumbering{Roman}               % 正文之前的页码用大写罗马字母编号.

\pagestyle{oldplain}\cleardoublepage
% !Mode:: "TeX:UTF-8"

%%% 说明: 此部分需要自己填写的内容:  (1) 中文摘要及关键词 (2) 英文摘要及关键词
%--------------------------------------------------------------------------
%%%页眉设置
%
%\pagestyle{fancyplain}
%\fancyhf{}  %清除以前对页眉页脚的设置
%\renewcommand{\headrulewidth}{0.5pt}
%\fancyfoot[C]{-\,\thepage\,-}
%\fancypagestyle{plain}
%{
%\fancyhead[CE]{武汉大学博士学位论文}
%\fancyhead[CO]{}
%\renewcommand{\headrulewidth}{0.5pt}
%\fancyfoot{}                                    % clear all footer fields
%\fancyfoot[C]{-\,\thepage\,-}
%}

%%%%%%%%%%%%%%%%%%%%%%%
%%% ------------ 中文摘要 ---------------%%%
%%%%%%%%%%%%%%%%%%%%%%%
\begin{cnabstract}
本文主要介绍和讨论了武汉大学博士毕业论文的~\LaTeX~模板.
指明了基本使用方法, 强调了公式排版的一些细节问题, 也指出了一些常见的排版错误.

本文主要介绍和讨论了武汉大学博士毕业论文的~\LaTeX~模板.
指明了基本使用方法, 强调了公式排版的一些细节问题, 也指出了一些常见的排版错误.

本文主要介绍和讨论了武汉大学博士毕业论文的~\LaTeX~模板.
指明了基本使用方法, 强调了公式排版的一些细节问题, 也指出了一些常见的排版错误.本文主要介绍和讨论了武汉大学博士毕业论文的~\LaTeX~模板.
指明了基本使用方法, 强调了公式排版的一些细节问题, 也指出了一些常见的排版错误.

本文主要介绍和讨论了武汉大学博士毕业论文的~\LaTeX~模板.
指明了基本使用方法, 强调了公式排版的一些细节问题, 也指出了一些常见的排版错误.

本文主要介绍和讨论了武汉大学博士毕业论文的~\LaTeX~模板.
指明了基本使用方法, 强调了公式排版的一些细节问题, 也指出了一些常见的排版错误.

本文主要介绍和讨论了武汉大学博士毕业论文的~\LaTeX~模板.
指明了基本使用方法, 强调了公式排版的一些细节问题, 也指出了一些常见的排版错误.

本文主要介绍和讨论了武汉大学博士毕业论文的~\LaTeX~模板.
指明了基本使用方法, 强调了公式排版的一些细节问题, 也指出了一些常见的排版错误.

本文主要介绍和讨论了武汉大学博士毕业论文的~\LaTeX~模板.
指明了基本使用方法, 强调了公式排版的一些细节问题, 也指出了一些常见的排版错误.本文主要介绍和讨论了武汉大学博士毕业论文的~\LaTeX~模板.
指明了基本使用方法, 强调了公式排版的一些细节问题, 也指出了一些常见的排版错误.


本文主要介绍和讨论了武汉大学博士毕业论文的~\LaTeX~模板.
指明了基本使用方法, 强调了公式排版的一些细节问题, 也指出了一些常见的排版错误.

本文主要介绍和讨论了武汉大学博士毕业论文的~\LaTeX~模板.
指明了基本使用方法, 强调了公式排版的一些细节问题, 也指出了一些常见的排版错误.

本文主要介绍和讨论了武汉大学博士毕业论文的~\LaTeX~模板.
指明了基本使用方法, 强调了公式排版的一些细节问题, 也指出了一些常见的排版错误.

本文主要介绍和讨论了武汉大学博士毕业论文的~\LaTeX~模板.
指明了基本使用方法, 强调了公式排版的一些细节问题, 也指出了一些常见的排版错误.本文主要介绍和讨论了武汉大学博士毕业论文的~\LaTeX~模板.
指明了基本使用方法, 强调了公式排版的一些细节问题, 也指出了一些常见的排版错误.



本文主要介绍和讨论了武汉大学博士毕业论文的~\LaTeX~模板.
指明了基本使用方法, 强调了公式排版的一些细节问题, 也指出了一些常见的排版错误.


本文主要介绍和讨论了武汉大学博士毕业论文的~\LaTeX~模板.
指明了基本使用方法, 强调了公式排版的一些细节问题, 也指出了一些常见的排版错误.本文主要介绍和讨论了武汉大学博士毕业论文的~\LaTeX~模板.
指明了基本使用方法, 强调了公式排版的一些细节问题, 也指出了一些常见的排版错误.本文主要介绍和讨论了武汉大学博士毕业论文的~\LaTeX~模板.
指明了基本使用方法, 强调了公式排版的一些细节问题, 也指出了一些常见的排版错误.


本文主要介绍和讨论了武汉大学博士毕业论文的~\LaTeX~模板.
指明了基本使用方法, 强调了公式排版的一些细节问题, 也指出了一些常见的排版错误.本文主要介绍和讨论了武汉大学博士毕业论文的~\LaTeX~模板.
指明了基本使用方法, 强调了公式排版的一些细节问题, 也指出了一些常见的排版错误.本文主要介绍和讨论了武汉大学博士毕业论文的~\LaTeX~模板.
指明了基本使用方法, 强调了公式排版的一些细节问题, 也指出了一些常见的排版错误.


本文主要介绍和讨论了武汉大学博士毕业论文的~\LaTeX~模板.
指明了基本使用方法, 强调了公式排版的一些细节问题, 也指出了一些常见的排版错误.本文主要介绍和讨论了武汉大学博士毕业论文的~\LaTeX~模板.
指明了基本使用方法, 强调了公式排版的一些细节问题, 也指出了一些常见的排版错误.本文主要介绍和讨论了武汉大学博士毕业论文的~\LaTeX~模板.
指明了基本使用方法, 强调了公式排版的一些细节问题, 也指出了一些常见的排版错误.


本文主要介绍和讨论了武汉大学博士毕业论文的~\LaTeX~模板.
指明了基本使用方法, 强调了公式排版的一些细节问题, 也指出了一些常见的排版错误.本文主要介绍和讨论了武汉大学博士毕业论文的~\LaTeX~模板.
指明了基本使用方法, 强调了公式排版的一些细节问题, 也指出了一些常见的排版错误.本文主要介绍和讨论了武汉大学博士毕业论文的~\LaTeX~模板.
指明了基本使用方法, 强调了公式排版的一些细节问题, 也指出了一些常见的排版错误.


本文主要介绍和讨论了武汉大学博士毕业论文的~\LaTeX~模板.
指明了基本使用方法, 强调了公式排版的一些细节问题, 也指出了一些常见的排版错误.本文主要介绍和讨论了武汉大学博士毕业论文的~\LaTeX~模板.
指明了基本使用方法, 强调了公式排版的一些细节问题, 也指出了一些常见的排版错误.本文主要介绍和讨论了武汉大学博士毕业论文的~\LaTeX~模板.
指明了基本使用方法, 强调了公式排版的一些细节问题, 也指出了一些常见的排版错误.



本文主要介绍和讨论了武汉大学博士毕业论文的~\LaTeX~模板.
指明了基本使用方法, 强调了公式排版的一些细节问题, 也指出了一些常见的排版错误.本文主要介绍和讨论了武汉大学博士毕业论文的~\LaTeX~模板.
指明了基本使用方法, 强调了公式排版的一些细节问题, 也指出了一些常见的排版错误.本文主要介绍和讨论了武汉大学博士毕业论文的~\LaTeX~模板.
指明了基本使用方法, 强调了公式排版的一些细节问题, 也指出了一些常见的排版错误.本文主要介绍和讨论了武汉大学博士毕业论文的~\LaTeX~模板.
指明了基本使用方法, 强调了公式排版的一些细节问题, 也指出了一些常见的排版错误.






\end{cnabstract}
\vspace{1em}\par\vfill


%%%--------- 关键词 -------- %%%
\cnkeywords{毕业论文, \LaTeX{}, 模板}

%%%%%%%%%%%%%%%%%%%%%%%


%%%%%%%%%%%%%%%%%%%%%%%
%%% ------------ 英文摘要 ---------------%%%
%%%%%%%%%%%%%%%%%%%%%%%

\begin{enabstract}
本文主要介绍和讨论了武汉大学博士毕业论文的~\LaTeX~模板.
指明了基本使用方法, 强调了公式排版的一些细节问题, 也指出了一些常见的排版错误.

本文主要介绍和讨论了武汉大学博士毕业论文的~\LaTeX~模板.
指明了基本使用方法, 强调了公式排版的一些细节问题, 也指出了一些常见的排版错误.

本文主要介绍和讨论了武汉大学博士毕业论文的~\LaTeX~模板.
指明了基本使用方法, 强调了公式排版的一些细节问题, 也指出了一些常见的排版错误.本文主要介绍和讨论了武汉大学博士毕业论文的~\LaTeX~模板.
指明了基本使用方法, 强调了公式排版的一些细节问题, 也指出了一些常见的排版错误.

本文主要介绍和讨论了武汉大学博士毕业论文的~\LaTeX~模板.
指明了基本使用方法, 强调了公式排版的一些细节问题, 也指出了一些常见的排版错误.

本文主要介绍和讨论了武汉大学博士毕业论文的~\LaTeX~模板.
指明了基本使用方法, 强调了公式排版的一些细节问题, 也指出了一些常见的排版错误.

本文主要介绍和讨论了武汉大学博士毕业论文的~\LaTeX~模板.
指明了基本使用方法, 强调了公式排版的一些细节问题, 也指出了一些常见的排版错误.

本文主要介绍和讨论了武汉大学博士毕业论文的~\LaTeX~模板.
指明了基本使用方法, 强调了公式排版的一些细节问题, 也指出了一些常见的排版错误.

本文主要介绍和讨论了武汉大学博士毕业论文的~\LaTeX~模板.
指明了基本使用方法, 强调了公式排版的一些细节问题, 也指出了一些常见的排版错误.本文主要介绍和讨论了武汉大学博士毕业论文的~\LaTeX~模板.
指明了基本使用方法, 强调了公式排版的一些细节问题, 也指出了一些常见的排版错误.


本文主要介绍和讨论了武汉大学博士毕业论文的~\LaTeX~模板.
指明了基本使用方法, 强调了公式排版的一些细节问题, 也指出了一些常见的排版错误.

本文主要介绍和讨论了武汉大学博士毕业论文的~\LaTeX~模板.
指明了基本使用方法, 强调了公式排版的一些细节问题, 也指出了一些常见的排版错误.

本文主要介绍和讨论了武汉大学博士毕业论文的~\LaTeX~模板.
指明了基本使用方法, 强调了公式排版的一些细节问题, 也指出了一些常见的排版错误.

本文主要介绍和讨论了武汉大学博士毕业论文的~\LaTeX~模板.
指明了基本使用方法, 强调了公式排版的一些细节问题, 也指出了一些常见的排版错误.本文主要介绍和讨论了武汉大学博士毕业论文的~\LaTeX~模板.
指明了基本使用方法, 强调了公式排版的一些细节问题, 也指出了一些常见的排版错误.



本文主要介绍和讨论了武汉大学博士毕业论文的~\LaTeX~模板.
指明了基本使用方法, 强调了公式排版的一些细节问题, 也指出了一些常见的排版错误.


本文主要介绍和讨论了武汉大学博士毕业论文的~\LaTeX~模板.
指明了基本使用方法, 强调了公式排版的一些细节问题, 也指出了一些常见的排版错误.本文主要介绍和讨论了武汉大学博士毕业论文的~\LaTeX~模板.
指明了基本使用方法, 强调了公式排版的一些细节问题, 也指出了一些常见的排版错误.本文主要介绍和讨论了武汉大学博士毕业论文的~\LaTeX~模板.
指明了基本使用方法, 强调了公式排版的一些细节问题, 也指出了一些常见的排版错误.


本文主要介绍和讨论了武汉大学博士毕业论文的~\LaTeX~模板.
指明了基本使用方法, 强调了公式排版的一些细节问题, 也指出了一些常见的排版错误.本文主要介绍和讨论了武汉大学博士毕业论文的~\LaTeX~模板.
指明了基本使用方法, 强调了公式排版的一些细节问题, 也指出了一些常见的排版错误.本文主要介绍和讨论了武汉大学博士毕业论文的~\LaTeX~模板.
指明了基本使用方法, 强调了公式排版的一些细节问题, 也指出了一些常见的排版错误.


本文主要介绍和讨论了武汉大学博士毕业论文的~\LaTeX~模板.
指明了基本使用方法, 强调了公式排版的一些细节问题, 也指出了一些常见的排版错误.本文主要介绍和讨论了武汉大学博士毕业论文的~\LaTeX~模板.
指明了基本使用方法, 强调了公式排版的一些细节问题, 也指出了一些常见的排版错误.本文主要介绍和讨论了武汉大学博士毕业论文的~\LaTeX~模板.
指明了基本使用方法, 强调了公式排版的一些细节问题, 也指出了一些常见的排版错误.


本文主要介绍和讨论了武汉大学博士毕业论文的~\LaTeX~模板.
指明了基本使用方法, 强调了公式排版的一些细节问题, 也指出了一些常见的排版错误.本文主要介绍和讨论了武汉大学博士毕业论文的~\LaTeX~模板.
指明了基本使用方法, 强调了公式排版的一些细节问题, 也指出了一些常见的排版错误.本文主要介绍和讨论了武汉大学博士毕业论文的~\LaTeX~模板.
指明了基本使用方法, 强调了公式排版的一些细节问题, 也指出了一些常见的排版错误.


本文主要介绍和讨论了武汉大学博士毕业论文的~\LaTeX~模板.
指明了基本使用方法, 强调了公式排版的一些细节问题, 也指出了一些常见的排版错误.本文主要介绍和讨论了武汉大学博士毕业论文的~\LaTeX~模板.
指明了基本使用方法, 强调了公式排版的一些细节问题, 也指出了一些常见的排版错误.本文主要介绍和讨论了武汉大学博士毕业论文的~\LaTeX~模板.
指明了基本使用方法, 强调了公式排版的一些细节问题, 也指出了一些常见的排版错误.



本文主要介绍和讨论了武汉大学博士毕业论文的~\LaTeX~模板.
指明了基本使用方法, 强调了公式排版的一些细节问题, 也指出了一些常见的排版错误.本文主要介绍和讨论了武汉大学博士毕业论文的~\LaTeX~模板.
指明了基本使用方法, 强调了公式排版的一些细节问题, 也指出了一些常见的排版错误.本文主要介绍和讨论了武汉大学博士毕业论文的~\LaTeX~模板.
指明了基本使用方法, 强调了公式排版的一些细节问题, 也指出了一些常见的排版错误.本文主要介绍和讨论了武汉大学博士毕业论文的~\LaTeX~模板.
指明了基本使用方法, 强调了公式排版的一些细节问题, 也指出了一些常见的排版错误.





\end{enabstract}
\vspace{1em}\par\vfill

%%%------ 英文关键词 ------- %%%
\enkeywords{\LaTeX{}, \dots}


      % 加入中英文摘要.

%---把目录加入到书签---%%%%%%%%%%%%%%
\pdfbookmark[0]{目录}{toc}%%%%%%%%%%%%
\pagestyle{oldplain}
\tableofcontents
\pagestyle{oldplain}\cleardoublepage
\newpage  \pagestyle{fancy} \fancyfancy
%------------------------------------------------------------------------------
% \include{metafile/figuresdes}      % 加入图检索.


%------------------------------------------------------------------------------
% \include{metafile/tabledes}      % 加入表检索.




%------------------------------------------------------------------------------



%------------------------------------------------------------------------------
\mainmatter %% 以下是正文
\baselineskip=20pt  % 正文行距为 20 磅
%%%%%%%%%%%%%%%%%%%%%%%%%%%%%%%%%%%%


%%%=====================================================================%%%
\chapter{绪论}

本文在 TeX Live 2015 软件环境下, 使用 XeLaTeX 引擎, 可以顺利编译.
TeX Live 的安装方法, 请移步查看: \url{http://aff.whu.edu.cn/huangzh/}.


引用格式请采用upcite命令行\upcite{AlquezarS96,AlquezarS97}
，而不是cite命令行\cite{AlquezarS96,AlquezarS97}。





\section{研究意义}


{\kaishu 以下文字参看武汉大学图书馆网页:
 \begin{center}
 \url{http://www.lib.whu.edu.cn/web/index.asp?obj_id=457&menu=h}.
 \end{center}
 提交论文之前, 请到该网址查看要求是否有新变化.
}



经研究生院与图书馆共同商议决定, 武汉大学研究生在学位论文答辩通过后,
采用以下方式提交学位论文：首先进行电子版论文网上提交, 经图书馆审核通过后, 进行纸本论文提交.
% \section*{前期准备}

\begin{itemize}
  \item[一、] 请下载《武汉大学学位论文使用授权协议书》(一式两份), 一份论文作者保存, 一份留学校存档.
    留学校存档的协议书事先用钢笔或中性笔填写后,  装订在提交给学校图书馆的纸本论文末页.\footnote{\heiti 特别说明:
    本文前面所附的《武汉大学学位论文使用授权协议书》, 图书馆要求放在最后;
    而研究生院是要求发在文档前面(在论文原创性声明之后). }

    \item[二、]涉密论文缴送到武汉大学档案馆, 由档案馆加盖公章后到图书馆办理相关离校手续.

    \item[三、]电子版论文要求
    

\begin{enumerate}[1.]

  \item 论文的电子文本应采用 Word 或 PDF 格式编辑(不加密).
  \item 电子版全文包括的内容及顺序应与纸本一致: 包括中英文封面、郑重声明、中英文摘要、目录、引言、正文(含图表)、参考文献及附录, 并放在一个文档中.
  \item 电子版文档中不能有空白页、标记、彩色字、乱码.
  \item 目录的页码一定要与正文的章节以及附后的内容相符合. 论文正文页码须从第一页起, 正文之前部分(不包括封面)用罗马字母(I, II, \dots)编页.
\end{enumerate}
\end{itemize}





\begin{figure}[!ht]
\centering
\includegraphics[width=0.5\columnwidth]{data/global_figures/whulogo.pdf}
\caption{
whulogo
}
\label{whulogo}
\end{figure}






进入图书馆主页(\url{http://www.lib.whu.edu.cn})点击``博硕士论文提交'', 进入论文提交系统.
具体步骤和注意事项请参见论文提交过程演示(PPT).
论文提交成功的 3 个工作日后, 可查收 Email 或在``已通过论文名单查询''中查询论文是否审核通过.
如果提交的论文不合格, 请按邮件要求修改后再次进入系统提交论文.



电子版论文提交审核通过的, 请提交 1 份纸本论文到相应的论文纸本缴送地,
末页须装订一份填写好的协议书.



















向武汉大学图书馆提交电子文档, 需单独编译文件: 在文档选项中添加 forlib.

原因: 图书馆要求提交的电子文档不能有空白页、彩色文字.

({\kaishu 这个要求使我在编制模板时遇到了一点问题: 这会导致电子版与纸质版的页码不一致. 博士论文每章的起始页默认在奇数页, 这会不可避免地出现空白页.})


\vfill

本文档下载更新地址: \url{http://aff.whu.edu.cn/huangzh/}. 使用之前, 请移步查看是否有更新.

问题反馈及建议, 请联系: huangzh@whu.edu.cn.



























  


%%%=====================================================================%%%

\chapter{第二章}


\begin{itemize}
  \item[一、] 请下载《武汉大学学位论文使用授权协议书》(一式两份), 一份论文作者保存, 一份留学校存档.
    留学校存档的协议书事先用钢笔或中性笔填写后,  装订在提交给学校图书馆的纸本论文末页.\footnote{\heiti 特别说明:
    本文前面所附的《武汉大学学位论文使用授权协议书》, 图书馆要求放在最后;
    而研究生院是要求发在文档前面(在论文原创性声明之后). }

    \item[二、]涉密论文缴送到武汉大学档案馆, 由档案馆加盖公章后到图书馆办理相关离校手续.

    \item[三、]电子版论文要求
    

\begin{enumerate}[1.]

  \item 论文的电子文本应采用 Word 或 PDF 格式编辑(不加密).
  \item 电子版全文包括的内容及顺序应与纸本一致: 包括中英文封面、郑重声明、中英文摘要、目录、引言、正文(含图表)、参考文献及附录, 并放在一个文档中.
  \item 电子版文档中不能有空白页、标记、彩色字、乱码.
  \item 目录的页码一定要与正文的章节以及附后的内容相符合. 论文正文页码须从第一页起, 正文之前部分(不包括封面)用罗马字母(I, II, \dots)编页.
\end{enumerate}
\end{itemize}





   
  

%%%=====================================================================%%%

%%%=====================================================================%%%
\chapter{总结与展望}


\section{全文工作总结}
{\kaishu 以下文字参看武汉大学图书馆网页:
 \begin{center}
 \url{http://www.lib.whu.edu.cn/web/index.asp?obj_id=457&menu=h}.
 \end{center}
 提交论文之前, 请到该网址查看要求是否有新变化.
}

    经研究生院与图书馆共同商议决定, 武汉大学研究生在学位论文答辩通过后,
    采用以下方式提交学位论文：首先进行电子版论文网上提交, 经图书馆审核通过后, 进行纸本论文提交.

% \section*{前期准备}

\begin{itemize}
  \item[一、] 请下载《武汉大学学位论文使用授权协议书》(一式两份), 一份论文作者保存, 一份留学校存档.
    留学校存档的协议书事先用钢笔或中性笔填写后,  装订在提交给学校图书馆的纸本论文末页.\footnote{\heiti 特别说明:
    本文前面所附的《武汉大学学位论文使用授权协议书》, 图书馆要求放在最后;
    而研究生院是要求发在文档前面(在论文原创性声明之后). }

    \item[二、]涉密论文缴送到武汉大学档案馆, 由档案馆加盖公章后到图书馆办理相关离校手续.

    \item[三、]电子版论文要求
    

\begin{enumerate}[1.]

  \item 论文的电子文本应采用 Word 或 PDF 格式编辑(不加密).
  \item 电子版全文包括的内容及顺序应与纸本一致: 包括中英文封面、郑重声明、中英文摘要、目录、引言、正文(含图表)、参考文献及附录, 并放在一个文档中.
  \item 电子版文档中不能有空白页、标记、彩色字、乱码.
  \item 目录的页码一定要与正文的章节以及附后的内容相符合. 论文正文页码须从第一页起, 正文之前部分(不包括封面)用罗马字母(I, II, \dots)编页.
\end{enumerate}
\end{itemize}








\section{未来研究展望}

    进入图书馆主页(\url{http://www.lib.whu.edu.cn})点击``博硕士论文提交'', 进入论文提交系统.
    具体步骤和注意事项请参见论文提交过程演示(PPT).
    论文提交成功的 3 个工作日后, 可查收 Email 或在``已通过论文名单查询''中查询论文是否审核通过.
    如果提交的论文不合格, 请按邮件要求修改后再次进入系统提交论文.



    电子版论文提交审核通过的, 请提交 1 份纸本论文到相应的论文纸本缴送地,
    末页须装订一份填写好的协议书.






   

%%%=====================================================================%%%




%%%------------附录. 不需要的话, 可以删除.--------
% \appendix

% \chapter{测试1}

% 测试

% \chapter{测试2}

% 测试

% \chapter{测试3}

% 测试




%%%%%%%%%%%%%%%%%=== 参考文献 ==%%%%%%%%%%%%%%%%%%%
\cleardoublepage\phantomsection

\addcontentsline{toc}{chapter}{参考文献}

{\zihao{5}
\bibliography{bibdata/bibtex}
% \bibliography{data/sec1/sec1}
}



\backmatter

% !Mode:: "TeX:UTF-8"

%%% 说明: 此部分需要自己填写的内容:  (1) 攻博(硕)期间科研成果 (2) 致谢

%%%%%%%%%%%%%%%%%%%%%%%
%%%------- 攻博(硕)期间科研成果 -------%%%
%%%%%%%%%%%%%%%%%%%%%%%
\reseachresult

{\zihao{5}
\begin{enumerate}[{[1]}]
\item   作者, 文章题目, 期刊名, 年份(期数): 起止页码
\item   作者, 书名, 版次, 出版地：出版单位, 年份, 起止页码
\end{enumerate}
}












%%%%%%%%%%%%%%%%%%%%%%%
%%% --------------- 致谢 ------------- - %%%
%%%%%%%%%%%%%%%%%%%%%%%
\acknowledgement



感谢你, 感谢他和她, 感谢大家.














%%%%%%%%%%%%%%%%%%%%%%%%%%%%%%%%%%%%%%%
%%%%%%%--判断是否需要空白页-----------------------------
  \iflib
  \else
  \newpage
  %\thispagestyle{empty}
  \cleardoublepage
  \fi
%%%%%%%-------------------------------------------------

 %%%致谢、作品列表 ．

\cleardoublepage
\end{document}

